\section{Primitive Geometry}
\subsection{Point}
In all of geometry, a point is the most basic geometrical primitive, it denotes a location relative to some space, it has no size
nor does it have any direction.

A point can be defined in eucledian space using the "\textbf{cartesian}" way of representing, which is commonly known as
\textbf{cartesian coordiantes}:
\begin{center}
    $P = \begin{pmatrix}a_1 & a_2 & a_3 & ... & a_n\end{pmatrix}$
\end{center}

where:
\begin{itemize}
    \item $P$: is the point itself :)
    \item $a_i$: is the magnitude of each basis vectors of this eucledian space.
    \item $n$: denotes the point residing in an $n$-th dimension eucledian space.
\end{itemize}

For example, in 3D cartesian coordinates, the point $\begin{pmatrix}1 & 2 & 3\end{pmatrix}$ 
(assuming we use $\hat{x}$, $\hat{y}$ and $\hat{z}$ as the basis vectors) would represent this:
\begin{enumerate}
    \item Begin in the origin of the 3D eucledian space: $\begin{pmatrix}0 & 0 & 0\end{pmatrix}$.
    \item Move $1$ unit in the $\hat{x}$ direction.
    \item Move $2$ units in the $\hat{y}$ direction.
    \item Move $3$ units in the $\hat{z}$ direction.
\end{enumerate}


\subsection{Line}
TODO

\subsection{Ray}
\begin{center}
    \includegraphics[width=13em, height=13em]{primitive_geometry_ray}
\end{center}

A \emph{ray} is defined as a directed segment that has a
\emph{origin}(a.k.a. having a position), and a \emph{direction}. It is mathematically defined as follows:

\begin{center}
    $P=O+\hat{D}t$
\end{center}

where:
\begin{itemize}
    \item $P$: is any point along the ray.
    \item $O$: is the origin of the ray, where it starts.
    \item $\hat{D}$: is the unit vector representing the direction of the ray.
    \item $t$: represents the percent / factor that "moves" the point along the ray.
\end{itemize}

If $t > 0$ then $P$ would be along the ray's direction $\hat{D}$.
If $t < 0$ then $P$ would be along the negative ray's direction $-\hat{D}$.
If $t = 0$ then $P$ would be exactly at the origin point $O$.

\subsection{Plane}
A plane can be defined in numerous ways.

\textbf{Distance-Normal form}:
\begin{center}
    \includegraphics[width=13em, height=13em]{primitive_geometry_distance_normal_plane}
\end{center}

\begin{center}
    $P \cdot \hat{n} - d = 0$
\end{center}

where:
\begin{itemize}
    \item $P$: is any point on the plane.
    \item $\hat{n}$: is the unit vector that represents the direction / facing of the plane.
    \item $d$: is the offset / distance of the plane from the origin along the direction $\hat{n}$ where it is facing.
\end{itemize}

The plane is defined here as all the points that resides at an equal level of length
given a unit vector pointing in the direction in which it's facing.

\sgap

\textbf{Point-Normal form}:
\begin{center}
    \includegraphics[width=13em, height=13em]{primitive_geometry_point_normal_plane}
\end{center}

\begin{center}
    $(P' - P) \cdot \hat{n} = 0$
\end{center}

where:
\begin{itemize}
    \item $P'$: is any point on the plane.
    \item $\hat{n}$: is the unit vector that represents the direction / facing of the plane.
    \item $P$: is an already known point on the plane.
\end{itemize}

Recall that dot product of two vectors gets the "parallelness" of two vectors,
if two vectors are perpendicular, their dot product should be 0, which makes sense for any point of a plane
if we draw a vector from point $P$ to point $P'$, the vector should lie on the plane.



\subsubsection{Ray-Plane Intersection for Distance-Normal Form}
We can replace the $P$ in plane's Distance-Norm formula with the Ray's formula with respect to points along the ray.

TODO: image here :3

\begin{center}
    $(O + \hat{D}t) \cdot \hat{n} - d = 0$
\end{center}

Since, $O,\hat{D}$ and $\hat{n}$ are known when calculating the intersection,
we solve for $t$:

\begin{center}
    $t = \frac{d - \hat{n} \cdot O}{\hat{n} \cdot \hat{D}}$
\end{center}

If $t \geq 0$ then the intersection is either at the beginning of the ray($t=0$) or
in front of the ray.
If $t < 0$ then the intersection is behind the ray's direction.

\sgap

If $\hat{n} \cdot \hat{D} = 0$ then we know that the ray is actually perpendicular with
the plane's surface direction. It either misses entirely, or hits infinite points.

\sgap

After figuring out a valid value for $t$, obtaining the point of intersection should be as simple as
plugging(substituting) the value $t$ in the original Ray formula:

\begin{center}
    $P_{intersect}=O+\hat{D}t_{found}$
\end{center}

$P_{intersect}$ should be the exact point of intersection on the plane!

\subsubsection{Converting between Point-Normal form and Distance-Normal form}
TODO: $d=O \cdot \hat{n}$ and $O = d \hat{n}$

\subsubsection{Ray-Plane Intersection for Point-Normal Form}
Much like the Distance-Normal form, we can also replace the $P$ in plane's Point-Norm formula with the Ray's formula with respect to points along the ray.

TODO: image here :3

\begin{center}
    $(O + \hat{D}t - P) \cdot \hat{n} = 0$
\end{center}

Solving for $t$ we get:
\begin{center}
    $t = \frac{(P - O) \cdot \hat{n}}{\hat{D} \cdot \hat{n}}$
\end{center}

The rest should be exactly the same as Distance-Normal form Ray-Plane intersection :3


\subsubsection{Mirroring a Point with regards to the Plane}
If you think about mirrors, they "imitate" everything from one side, just "mirrored" / "flipped" using the mirror itself
as a center for flipping. We can think of a plane as a mirror in 3D space, we can mirror any point from one side to the other side
of the "mirror"(plane) using some formula :O

\sgap
The formula for mirroring a point in \textbf{Point-Normal form} for the plane should be as follows:
\begin{center}
    $P' = P - 2 \hat{n} [(P - O) \cdot \hat{n}]$
\end{center}

where:
\begin{itemize}
    \item $P$: is the initial point to mirror.
    \item $P'$: is the final point that got mirrored.
    \item $\hat{n}$: is the normal / direction that the plane is facing.
    \item $O$: is any known point on the plane.
\end{itemize}

This is how it looks like visually:

TODO: put an image here showing the plane, the point $P$ and the reflected point $P'$, and also the projected vector
$\hat{n}[(P - O) \cdot \hat{n}]$





\subsubsection{Plane from two different, non-zero vectors}
You can also define a plane's facing direction with two different, non-zero vectors, in mathematical terms:

TODO: image here :3

\begin{center}
    $\hat{n} = \frac{\vec{a} \times \vec{b}}{\|\vec{a} \times \vec{b}\|}$
\end{center}

where $\vec{a} \neq \vec{b}, \|\vec{a}\| \neq 0, \|\vec{b}\| \neq 0$




\subsection{Sphere}
A sphere is defined as all the points $P$ at an equal distance $r$
from the center point $C$:

\begin{center}
    $\|P-C\|-r=0$
\end{center}

where:
\begin{itemize}
    \item $P$: is any point on the surface of the sphere.
    \item $C$: is the center of the sphere.
    \item $r$: is the radius of the circle, from the center $C$ to any point on the
        surface of the sphere $P$.
\end{itemize}

\subsubsection{Ray-Sphere Intersection}
TODO

\subsection{Axis-Aligned Bounding Box}
TODO

\subsection{Triangle}
TODO

\subsection{Barycentric Coordinates}
TODO

